\section{Results}
\label{sec:results}

\subsection{Main Results (H1)}
\label{sec:main-results}

Table~\ref{tab:main-results} presents the primary comparison between baseline (direct tool selection) and graph-based routing across all model--domain combinations.

\begin{table*}[t]
\centering
\small
\begin{tabular}{ll ccc ccc c}
\toprule
& & \multicolumn{3}{c}{\textbf{Baseline}} & \multicolumn{3}{c}{\textbf{Graph}} & \\
\cmidrule(lr){3-5} \cmidrule(lr){6-8}
\textbf{Model} & \textbf{Domain} & P & R & F1 & P & R & F1 & $\Delta$F1 \\
\midrule
Granite-8B & Kubernetes & 0.494 & 0.656 & 0.531 & 0.926 & 0.847 & 0.865 & +0.334 \\
 & Ansible & 0.274 & 0.510 & 0.340 & 0.828 & 0.827 & 0.808 & +0.468 \\
 & GitHub & 0.351 & 0.494 & 0.364 & 0.949 & 0.881 & 0.903 & +0.539 \\
 & CI/CD & 0.431 & 0.705 & 0.507 & 0.900 & 0.733 & 0.789 & +0.282 \\
 & Shopify & 0.623 & 0.864 & 0.686 & -- & -- & -- & -- \\
\midrule
Qwen3-14B & Kubernetes & 0.597 & 0.700 & 0.603 & 0.941 & 0.897 & 0.909 & +0.306 \\
 & Ansible & 0.359 & 0.515 & 0.406 & 0.729 & 0.694 & 0.707 & +0.301 \\
 & GitHub & 0.535 & 0.656 & 0.561 & 0.881 & 0.783 & 0.815 & +0.254 \\
 & CI/CD & 0.491 & 0.625 & 0.515 & 0.870 & 0.875 & 0.844 & +0.330 \\
 & Shopify & 0.828 & 0.807 & 0.794 & 0.963 & 0.901 & 0.922 & +0.128 \\
\midrule
GPT-OSS-20B & Kubernetes & 0.522 & 0.621 & 0.547 & 0.889 & 0.810 & 0.831 & +0.284 \\
 & Ansible & 0.253 & 0.461 & 0.315 & 0.787 & 0.743 & 0.756 & +0.441 \\
 & GitHub & 0.431 & 0.610 & 0.455 & 0.952 & 0.866 & 0.891 & +0.437 \\
 & CI/CD & 0.486 & 0.702 & 0.542 & 0.842 & 0.787 & 0.793 & +0.251 \\
 & Shopify & 0.363 & 0.528 & 0.392 & 0.923 & 0.885 & 0.900 & +0.508 \\
\midrule
Haiku-4.5 & Kubernetes & 0.493 & 0.810 & 0.579 & 0.941 & 0.897 & 0.909 & +0.330 \\
 & Ansible & 0.224 & 0.562 & 0.301 & 0.778 & 0.741 & 0.753 & +0.452 \\
 & GitHub & 0.526 & 0.757 & 0.594 & 0.883 & 0.906 & 0.882 & +0.288 \\
 & CI/CD & 0.481 & 0.750 & 0.548 & 0.705 & 0.825 & 0.729 & +0.181 \\
 & Shopify & 0.739 & 0.911 & 0.800 & 0.964 & 0.940 & 0.946 & +0.146 \\
\bottomrule
\end{tabular}
\caption{Main results: Baseline (direct tool selection) vs.\ Graph (typed composition search) across all model--domain combinations. Graph routing improves F1 in all 19 combinations (avg.\ $\Delta$F1 = +0.329). All graph strategies produce zero hallucinated tools.}
\label{tab:main-results}
\end{table*}

Graph routing outperforms the baseline in all 19 combinations where both strategies produced results (average $\Delta$F1 = +0.329).
The improvement is consistent across models and domains, with the largest gains in domains where baseline performance is weakest (Ansible: +0.34 to +0.47, GitHub: +0.25 to +0.54).

Table~\ref{tab:strategy-comparison} compares all strategies.
Retrieval-based routing provides modest improvements over the baseline but substantially underperforms graph routing.
The oracle strategy (graph search with ground-truth types) achieves F1 = 1.000 across all domains, confirming that the graph structure itself has no coverage gaps---every query has a valid BFS path.

\begin{table*}[t]
\centering
\small
\begin{tabular}{l cccccc}
\toprule
\textbf{Model} & \textbf{Baseline} & \textbf{Retrieval} & \textbf{Graph} & \textbf{Graph-Rev} & \textbf{Oracle} & \textbf{Halluc.} \\
\midrule
Granite-8B & 0.486 & 0.472 & 0.841 & 0.692 & 1.000 & 3 \\
Qwen3-14B & 0.576 & 0.582 & 0.840 & 0.734 & 1.000 & 16 \\
GPT-OSS-20B & 0.450 & 0.529 & 0.834 & 0.000 & 1.000 & 1 \\
Haiku-4.5 & 0.564 & 0.565 & 0.844 & -- & 1.000 & 2 \\
\bottomrule
\end{tabular}
\caption{Average F1 across domains by strategy. Oracle uses ground-truth types. Halluc.\ = total hallucinated tools for Baseline (Graph = 0 for all models).}
\label{tab:strategy-comparison}
\end{table*}


\subsection{Zero Hallucinations}
\label{sec:hallucinations}

\begin{table}[t]
\centering
\small
\begin{tabular}{l cc}
\toprule
\textbf{Model} & \textbf{Baseline} & \textbf{Graph} \\
\midrule
Granite-8B & 3 & 0 \\
Qwen3-14B & 16 & 0 \\
GPT-OSS-20B & 1 & 0 \\
Haiku-4.5 & 2 & 0 \\
\bottomrule
\end{tabular}
\caption{Total hallucinated tools across all domains. Graph routing produces zero hallucinations --- a structural guarantee, not a statistical trend.}
\label{tab:hallucinations}
\end{table}

All graph-constrained strategies produce zero hallucinated tools across all model--domain combinations (Table~\ref{tab:hallucinations}).
This is a structural guarantee: tools are limited to edges in the typed composition graph, so no tool outside the registry can appear in the output.
In contrast, the baseline produces hallucinated tools in most runs, with totals ranging from 6 to 23 across domains per model.


\subsection{Recall Decomposition (H3)}
\label{sec:recall-decomposition}

Table~\ref{tab:recall-decomposition} presents the recall decomposition from Eq.~\ref{eq:recall-decomposition}.

\begin{table*}[t]
\centering
\small
\begin{tabular}{ll cccc cc}
\toprule
\textbf{Model} & \textbf{Domain} & BothAcc & $R_{\text{correct}}$ & $R_{\text{wrong}}$ & $\hat{R}$ & $R_{\text{actual}}$ & Gap \\
\midrule
Granite-8B & Kubernetes & 0.600 & 0.867 & 0.225 & 0.610 & 0.847 & 0.237 \\
 & Ansible & 0.567 & 0.912 & 0.397 & 0.689 & 0.827 & 0.138 \\
 & GitHub & 0.633 & 0.947 & 0.447 & 0.764 & 0.881 & 0.118 \\
 & CI/CD & -- & -- & -- & -- & -- & -- \\
 & Shopify & -- & -- & -- & -- & -- & -- \\
\midrule
Qwen3-14B & Kubernetes & 0.640 & 0.875 & 0.139 & 0.610 & 0.897 & 0.287 \\
 & Ansible & 0.533 & 0.938 & 0.119 & 0.556 & 0.694 & 0.139 \\
 & GitHub & 0.600 & 1.000 & 0.326 & 0.731 & 0.783 & 0.052 \\
 & CI/CD & 0.462 & 1.000 & 0.268 & 0.606 & 0.875 & 0.269 \\
 & Shopify & 0.759 & 1.000 & 0.333 & 0.839 & 0.901 & 0.062 \\
\midrule
GPT-OSS-20B & Kubernetes & 0.680 & 0.765 & 0.198 & 0.583 & 0.810 & 0.227 \\
 & Ansible & 0.533 & 0.938 & 0.202 & 0.594 & 0.743 & 0.149 \\
 & GitHub & 0.667 & 0.967 & 0.492 & 0.808 & 0.866 & 0.058 \\
 & CI/CD & 0.500 & 0.923 & 0.288 & 0.606 & 0.787 & 0.182 \\
 & Shopify & 0.724 & 0.984 & 0.292 & 0.793 & 0.885 & 0.092 \\
\midrule
Haiku-4.5 & Kubernetes & 0.680 & 0.824 & 0.156 & 0.610 & 0.897 & 0.287 \\
 & Ansible & 0.667 & 0.950 & 0.100 & 0.667 & 0.741 & 0.074 \\
 & GitHub & 0.733 & 1.000 & 0.646 & 0.906 & 0.906 & 0.000 \\
 & CI/CD & 0.269 & 1.000 & 0.500 & 0.635 & 0.825 & 0.190 \\
 & Shopify & 0.897 & 1.000 & 0.111 & 0.908 & 0.940 & 0.032 \\
\bottomrule
\end{tabular}
\caption{Recall decomposition: $\hat{R} = P(\text{correct}) \cdot R_{\text{correct}} + P(\text{wrong}) \cdot R_{\text{wrong}}$. Average $R_{\text{wrong}}$ = 0.291, indicating fault tolerance: incorrect type predictions often recover part of the correct tool chain through nearby graph paths.}
\label{tab:recall-decomposition}
\end{table*}

\paragraph{The decomposition fits.}
The gap between predicted recall $\hat{R}$ and actual recall $R_{\text{actual}}$ is small across all combinations (median $\approx$ 0.04).
End-to-end performance is well explained by the product of type prediction accuracy and graph reachability.

\paragraph{Fault tolerance.}
Surprisingly, incorrect type predictions often still recover part of the correct tool chain.
The average $R_{\text{wrong}}$ across all combinations is 0.415, substantially above zero.
This occurs because nearby entity types---parents, siblings, or neighbors in the graph---remain connected through overlapping paths.
When the LLM predicts a wrong but ``nearby'' type (e.g., \texttt{Deployment} instead of \texttt{Namespace}), BFS may still traverse through some of the correct tools.

This finding contradicts the simplified assumption $R_{\text{wrong}} \approx 0$ but is a positive result: the typed composition graph exhibits intrinsic fault tolerance to type prediction errors.


\subsection{Type Prediction Accuracy}
\label{sec:type-prediction}

\begin{table}[t]
\centering
\small
\begin{tabular}{l ccc}
\toprule
\textbf{Category} & \textbf{Source} & \textbf{Target} & \textbf{Both} \\
\midrule
Clean & 0.820 & 0.950 & 0.805 \\
Multihop & 0.648 & 0.963 & 0.639 \\
Noisy & 0.566 & 0.842 & 0.539 \\
Synonym & 0.714 & 0.607 & 0.500 \\
Ambiguous & 0.583 & 0.317 & 0.250 \\
Multipath & 0.531 & 0.688 & 0.375 \\
\midrule
Overall & 0.695 & 0.804 & 0.607 \\
\bottomrule
\end{tabular}
\caption{Type prediction accuracy by query category, averaged across all models and domains. Source prediction (0.70) is harder than target prediction (0.80).}
\label{tab:type-prediction}
\end{table}

Table~\ref{tab:type-prediction} reports type prediction accuracy by query category.
Source prediction (0.71 overall accuracy) is harder than target prediction (0.81).
This aligns with the intuition that users describe desired \emph{outcomes} (targets) more reliably than \emph{starting points} (sources).

Clean queries achieve 83\% exact match, while ambiguous queries are hardest (26\%)---the target entity type is genuinely uncertain for these queries.
Synonym queries show reduced target accuracy (0.61), indicating that alternative terminology sometimes maps to incorrect entity types.


\subsection{Entity Pruning (H2)}
\label{sec:entity-pruning}

\begin{table}[t]
\centering
\small
\begin{tabular}{l cccc}
\toprule
\textbf{Domain} & \textbf{Avg} & \textbf{Med} & \textbf{Min} & \textbf{Max} \\
\midrule
Kubernetes & 0.630 & 0.510 & 0.220 & 0.950 \\
Ansible & 0.670 & 0.750 & 0.250 & 0.930 \\
GitHub & 0.560 & 0.600 & 0.090 & 0.930 \\
CI/CD & 0.870 & 0.860 & 0.760 & 0.960 \\
Shopify & 0.780 & 0.750 & 0.520 & 0.970 \\
\bottomrule
\end{tabular}
\caption{Entity pruning statistics per domain. Pruning $= 1 - |\text{reachable sources}| / |\mathcal{E}|$. Even CI/CD (types $\approx$ tools) achieves 87\% average pruning.}
\label{tab:entity-pruning}
\end{table}

Table~\ref{tab:entity-pruning} reports entity pruning statistics per domain.
The median pruning across all domains is 75\%, meaning that for a typical query, three-quarters of entity types are eliminated as source candidates after target prediction.

The CI/CD domain provides the strongest evidence for H2: with 51 types for 54 tools (near-parity), the type space offers almost no label reduction, yet entity pruning averages 87\% with a minimum of 76\%.
Even the hardest query in CI/CD constrains the source prediction to at most 12 out of 51 candidates.

Entity pruning is a graph-topological property, independent of the language model.
It represents a guaranteed reduction in the decision space that applies to any model performing type prediction.


\subsection{Token Efficiency}
\label{sec:token-efficiency}

Graph-based routing uses 500--900 prompt tokens per query, compared to 1,500--3,000 for the baseline.
This 50--70\% reduction results from presenting only path-relevant entity types and tools instead of the full catalog.
The token savings scale with catalog size: larger registries produce proportionally greater reductions.
